\section{Finite-type clearance on the actual roof fibers}
\label{sec:v62-finite-type-clearance}
The determinant cutoff of Section~\ref{sec:v61-clearance-coarea} leaves
points at which the clearance is tangent to a regular roof fiber. A
higher derivative along that fiber can still give a positive height
gain. We prove this on the original physical words, including their
exact labels. The estimate applies to the finite-type part of the
first-clearance source; no finite-type assertion is imposed on its
complement.

\subsection{A one-dimensional estimate}
\begin{lemma}[A derivative sublevel bound]\label{lem:v62-sublevel}
For each integer $r\ge1$ there is $C_r$ with the following property.
If $V\subset\R$ is an interval, $f\in C^r(V)$, and
$|f^{(r)}|\ge b>0$ on $V$, then, for an interval $H$ of length $s$,
\begin{equation}\label{eq:v62-sublevel}
 |\{v\in V:f(v)\in H\}|\le C_r(s/b)^{1/r}.
\end{equation}
The interval $V$ need not include its endpoints.
\end{lemma}
\begin{proof}
Exhaust by compact subintervals. Let $E$ be the sublevel set in one
such interval and write $l=|E|$. If $l>0$, a greedy selection gives
$r+1$ points of $E$ separated by at least $l/(4(r+1))$: the union of
$r$ neighborhoods of this radius has length less than $l$. Write
them in increasing order as $v_0,\ldots,v_r$, and subtract one fixed
endpoint of $H$ from $f$. The resulting values lie in $[0,s]$.
The Lagrange formula for their divided difference bounds its absolute
value by
$(r+1)s(4(r+1)/l)^r$. Repeated Rolle's theorem identifies this
difference with $f^{(r)}(\xi)/r!$ for some $\xi$ between the extreme
points. Hence
$b\le r!(r+1)s(4(r+1)/l)^r$.
This proves the estimate with a constant depending only on $r$.
Approximation of $E$ from within and monotone convergence remove the
compact exhaustion. For $s=0$, apply the result to decreasing positive
lengths.
\end{proof}

\subsection{The derivative on a physical fiber}
Retain the first-clearance witness $(j,d)$ and its ordering from
Section~\ref{sec:v61-clearance-coarea}. Recenter at the actual image
collision $y=T_R^{j+1}x$. In a fixed rational angular chart write
\[
 \alpha=\alpha_*+2\arctan x_1,\qquad p=x_2,\qquad |x_1|<1.
\]
Use a fixed finite cover and assign overlaps in a fixed order when a
single source partition is needed. The invariant density in these
coordinates is at most $1/(2\pi)$. On a regular, actual two-sided word
through $y$ put
\begin{equation}\label{eq:v62-fiber-objects}
 F(y)=L_{m,R}(T_R^{-(j+1)}y),\qquad Z(y)=Z_d(y).
\end{equation}
The function $F$ sums the original $m$ flights, including those before
the mark. It is not a roof for a nearby continued word. For $i=1,2$,
write $i'=3-i$ and, where $F_i=\partial_iF\ne0$, set
\begin{equation}\label{eq:v62-fiber-derivative}
 \mathscr D_i=\partial_{i'}-\frac{F_{i'}}{F_i}\partial_i.
\end{equation}
On a graph of $F=t$, parametrized by $x_{i'}$, this is its ordinary
derivative. In $\mathscr D_i^rZ$ all coefficients of the vector field
are differentiated as well.

For $a,b>0$, define the order-$r$ carrier at the selected witness by
\begin{equation}\label{eq:v62-type-carrier}
 \mathcal T_r(a,b)=
 \bigcup_{\text{assigned charts},\ i=1,2}
       \{|F_i|\ge a,\ |\mathscr D_i^rZ|\ge b\}.
\end{equation}
All functions in this definition are evaluated on the actual word of
the source point. This definition includes nonzero first derivatives
when $r=1$. When $r=2$, it also permits $\det D(F,Z)=0$ provided the
second derivative and the chosen roof derivative have the stated
lower bounds.

\begin{lemma}[Finite-order physical fiber complexity]
\label{lem:v62-jet-complexity}
For each fixed $r\ge1$ there are $C_r>0$, $A_r>1$ such that the
parts of $F=t$ satisfying either pivot condition in
\eqref{eq:v62-type-carrier} can be covered, apart from finitely many
points, by at most $C_rA_r^m$ graph intervals on which that condition
holds throughout. The bound includes all original center words and
charts and is independent of $R,t,a,b,j,d$.
\end{lemma}
\begin{proof}
We use the auxiliary physical graph in
Lemma~\ref{lem:v44-level-complexity}, without eliminating intermediate
collisions. The initial data are now at the mark. Forward and backward
flight equations determine each of the $m$ contacts and velocities;
all selected incidence cosines are positive on a regular word.
For each flight, implicit differentiation of its quadratic contact
equation determines the first jet uniquely. Its coefficient for the
flight-time derivative is the nonzero incoming normal component.
The reflection equation is polynomial in the normal and velocity.
Differentiating these equations through order $r$ determines each
successive jet linearly in its highest-order unknown, using already
determined lower jets. The same construction applies to the reversed
flights. It introduces at most $C_r(m+1)$ variables and equations of
degree bounded solely in terms of $r$. Reciprocal variables with their
nonzero sign conditions record denominators; no inverse is taken at
a grazing contact.

Sum the time jets to obtain the jets of $F$, and use the scalar-product
formula for $Z_d$ to obtain its jets. Add $F=t$, the pivot inequalities,
and the rational expression for $\mathscr D_i^rZ$ with $|F_i|\ge a$.
The number of polynomial sign conditions is still at most $C_r(m+1)$
and their degrees remain bounded in terms of $r$. The auxiliary jets
are unique, so their graph projects homeomorphically onto the regular
source set. The component estimate \eqref{eq:v44-sign-bound}, namely
\[
 d(2d-1)^{v-1}\sum_{h=0}^{v}4^h\binom{s}{h}
       \le d(2d-1)^{v-1}5^s,
\]
from \cite[Section 3.2, equation (3.3)]{BPR}, is therefore exponential
in $m$, not exponential in $m\log m$. Here $v,s=O_r(m+1)$ and $d$
is fixed after $r$. The bound is independent of polynomial coefficients,
so the thresholds and roof value do not change it.

To make the graph-interval conclusion explicit, a one-dimensional
component in the pivot region is a regular level curve. Its tangent
can be oriented to have positive $x_{i'}$ component, because $F_i$
has a fixed nonzero sign. The projection to $x_{i'}$ is consequently
strictly increasing along that oriented component. It has no closed
component or turning point. Subdivision at endpoints and at the
finitely many zero-dimensional strata gives graphs over intervals.
Their number is bounded by applying the same component estimate to
the sign strata and their frontier conditions. Equivalently these
frontiers are obtained by adjoining equality in the finitely many
pivot, jet and physical inequalities already present. A one-dimensional
regular level has no branching in the pivot region. This subdivision
thus keeps the exponential bound. The number of candidate center
words is itself at most $D^m$, and the chart cover is fixed.

Only the regular physical graph and its interval pieces are used.
A physical seam may end a piece; it is not filled in by a continuation
of the orbit. The section and first-defect indicators can restrict a
piece further, but are not needed in this positive upper cover.
\end{proof}

\begin{theorem}[Positive finite-type clearance height]
\label{thm:v62-finite-type-height}
Restrict the original positive first-clearance source to
$\mathcal T_r(a,b)$, and denote its exact roof density with insertion
$w$ by $b^{\varepsilon,\mathrm{ft}_r(a,b),w}_{n,k,m,R}$.
For each fixed $r\ge1$ there are $C_r>0$, $A_r>1$ such that
\begin{equation}\label{eq:v62-finite-type-height}
 \operatorname*{ess\,sup}_t\sum_{n=1}^m\sum_{k\in\Z^2}
 |b^{\varepsilon,\mathrm{ft}_r(a,b),w}_{n,k,m,R}(t)|
 \le C_r M A_r^m a^{-1}(\varepsilon/b)^{1/r}
\end{equation}
for $m\ge2$, $0<2\varepsilon<s_0$, $a,b>0$ and $|w|\le M$.
The thresholds may depend on $m$; the constants do not. The source
outside the carrier is not discarded or assigned this estimate.
\end{theorem}
\begin{proof}
Change the source origin to $T_R^{j+1}x$. Invariance keeps its density
bounded by $1/c_*$ times the collision density, even with the original
section and exact-return conditions retained. On one pivot graph,
let $X_i(t,v)$ have its $i'$th coordinate equal to $v$ and its
$i$th coordinate given by the implicit graph. Ordinary coarea gives
\[
 \int_{F=t}\frac{r(y)}{|\nabla F|}\,d\mathcal H^1(y)
   =\int\frac{r(X_i(t,v))}{|F_i(X_i(t,v))|}\,dv,
       \qquad 0\le r\le\frac1{2\pi c_*}.
\]
It may first be applied on compact regular subsets and then exhausted
by positivity. At the selected witness the original clearance source
is supported in the two intervals
$Z\in[R,R+s_j]\cup[-R-s_j,-R]$, with
$s_j=2\varepsilon q_0^{\min(j,m-1-j)/4}$.
Lemma~\ref{lem:v62-sublevel}, applied along each graph in
Lemma~\ref{lem:v62-jet-complexity}, bounds the length contributing to
this integral by $C_r(s_j/b)^{1/r}$. The pivot denominator costs $a^{-1}$.
Sum over the two pivots, the fixed charts and candidate disks, and all
physical words. Finally
\[
 \sum_{j=0}^{m-1}s_j^{1/r}
       \le \frac{2(2\varepsilon)^{1/r}}{1-q_0^{1/(4r)}}.
\]
This proves the positive estimate. For fixed $m$ the events
$\Pi_{n,k,m,R}$ are disjoint. Domination before pushforward therefore
proves the sum of absolute inserted densities without a factor for
the number of targets. Only a Lebesgue-null set of roofs is exceptional.
\end{proof}

\begin{corollary}[Fold clearance and the retained positive source]
\label{cor:v62-fold-clearance}
On $\mathcal T_2(a,b)$ the bound is
$CMA_2^m a^{-1}\sqrt{\varepsilon/b}$, including points with vanishing
first clearance derivative. For $1\le r\le r_*$ and thresholds $b_r$,
assign each point of $\bigcup_{r\le r_*}\mathcal T_r(a,b_r)$ to its
first qualifying order, and retain the rest as
$b^{\varepsilon,\mathrm{rest},1}$. Then
\[
 b^{\varepsilon,\mathrm{clr},1}
    =\sum_{r=1}^{r_*}b^{\varepsilon,[r],1}
                         +b^{\varepsilon,\mathrm{rest},1},
       \qquad b^{\varepsilon,[r],1},b^{\varepsilon,\mathrm{rest},1}\ge0.
\]
With $P,G,e_{B,m}$ as in Section~\ref{sec:v61-clearance-coarea},
$\varepsilon=\varepsilon(B)$, and supremum over all $R,n,k$,
\begin{align}\label{eq:v62-finite-type-error}
 \|P-G-m^2b^{\varepsilon,\mathrm{rest},1}\|_\infty
 \le e_{B,m}+Cm^2A^m\varepsilon
       +m^2\sum_{r=1}^{r_*}C_rA_r^m a^{-1}(\varepsilon/b_r)^{1/r}.
\end{align}
\end{corollary}
\begin{proof}
Use the positive partition, Theorem~\ref{thm:v62-finite-type-height},
and the first-incidence bound of
Corollary~\ref{cor:v61-incidence-height-gain} in the exact positive
raw-error representation. The Gaussian reference $G$ is unchanged,
including its finite arithmetic factor and zero classes.
\end{proof}

The estimate covers physical finite-type strata rather than assuming
that every clearance point has finite type with a uniform constant.
In particular, a critical point of the total roof need not lie in any
pivot carrier. The factors $A_r^m$ also remain explicit. Thus the
corollary is a positive source reduction at finite collision count;
it is not the ordered central-scale clearance-height limit.
